\Needspace{8\baselineskip}\originalchapter{2002 年碰撞与流体作业}
\originalsection{2002 PS3：自由程、聚变 α 粒子与逃逸电子}
\begin{exercise}\label{ex:hist02-3-1}
PS3 原题1（有官方答案）。按特征速度除动量损失频率定义平均自由程，计算：(a) $T_e=T_i=4\,$keV、$n=10^{14}\,\mathrm{cm^{-3}}$ 氢 tokamak 的热电子；(b) 同一电子与质子的动量交换自由程 $\lambda_c$ 与 Debye 长度之比，推导它与基本等离子体参数的关系；(c) 同一热离子；(d) $T_e=5\,$eV、$n=5\times10^{12}\,\mathrm{cm^{-3}}$ 的加工等离子体热电子。
\end{exercise}
\begin{solution}
编者独立解答。采用 $v_{ts}=\sqrt{2T_s/m_s}$。氢等离子体的常用热动量频率近似为
\begin{equation}\nu_{ei}=2.91\times10^{-6}\frac{n_{\rm cm^{-3}}\ln\Lambda}{T_{e,\rm eV}^{3/2}},\qquad \nu_{ii}=4.80\times10^{-8}\frac{n_{\rm cm^{-3}}\ln\Lambda}{T_{i,\rm eV}^{3/2}}\quad(\mathrm{s^{-1}}).\end{equation}
系数依“热频率”定义有量级一差。以量子和经典近距较大者为截止，前一种环境可取 $\ln\Lambda\simeq17$，后一种约 $11$。代入给
\begin{equation}\lambda_e(4\mathrm{keV})\simeq1.92\times10^3\,\mathrm m,\qquad
\lambda_i(4\mathrm{keV})\simeq2.71\times10^3\,\mathrm m,\qquad
\lambda_e(5\mathrm{eV})\simeq9.2\times10^{-2}\,\mathrm m.
\end{equation}
首环境 $\lambda_D=\sqrt{\epsilon_0T_e/(ne^2)}\simeq4.70\times10^{-5}\,$m，故 $\lambda_e/\lambda_D\sim4.1\times10^7$。由 $\nu\sim ne^4\ln\Lambda/(\epsilon_0^2m^{1/2}T^{3/2})$ 和 $\lambda_D^2=\epsilon_0T/(ne^2)$，可见
\begin{equation}\frac{\lambda_c}{\lambda_D}\sim\frac{4\pi n\lambda_D^3}{\ln\Lambda}\equiv\frac{N_D}{\ln\Lambda},\end{equation}
其中系数依 Debye 球、碰撞率定义改变，但弱耦合等离子体的 $N_D\gg1$ 决定比例很大。长自由程指该局域二体碰撞模型下弱碰撞，不意味装置不存在几何损失。
\end{solution}

\begin{exercise}\label{ex:hist02-3-2}
PS3 原题2（有官方答案）。50:50 D--T 等离子体 $n\sim10^{14}\,\mathrm{cm^{-3}}$、$T_e=T_D=T_T=10\,$keV 中产生 $3.5\,$MeV 的 $\alpha$ 粒子。估计与电子、氘、氚三物种碰撞的能量松弛时间，指出开始减速的主要物种及减速过程中能量最终主要加热谁。须区分快 $\alpha$ 与接近热速阶段。
\end{exercise}
\begin{solution}
编者独立解答。初速 $v_\alpha=\sqrt{2E_\alpha/m_\alpha}\simeq1.30\times10^7\,$m/s；$10\,$keV 电子热速约 $5.93\times10^7\,$m/s，而 D、T 离子热速约 $9.8,8.0\times10^5\,$m/s。因而开始时 $\alpha$ 对离子是快试探粒子，对电子却在电子热分布的低速区。对背景 $b$ 可用 Landau 阻力函数
\begin{equation}
\nu_{s,b}(v)=\frac{n_bZ_\alpha^2Z_b^2e^4\ln\Lambda}{4\pi\epsilon_0^2m_\alpha m_bv^3}\left[\operatorname{erf}(x_b)-\frac{2x_b}{\sqrt\pi}e^{-x_b^2}\right],\quad x_b=\frac{v}{\sqrt{2T_b/m_b}},
\end{equation}
并明确不同教科书对 $\nu_s$ 前因子的定义可能不同。代入 $n_D=n_T=n_e/2$ 后，低 $x_e$ 有括号 $\sim4x_e^3/(3\sqrt\pi)$，电子拖曳率近似常数；离子 $x_i\gg1$ 有括号近 1、拖曳率随 $v^{-3}$ 降。由两率相等定义临界速度 $v_c$，于是 $E_\alpha\gg E_c\sim\frac12m_\alpha v_c^2$ 的首段主要向电子给能，减慢至 $E\lesssim E_c$ 后离子分走越来越多能量。$E_c$ 的量级为数百 keV，且 D/T 的贡献须分别以 $n_b/m_b$ 权重求和；以 $\ln\Lambda=16$ 按本题官方扫描所用 NRL 频率规范独立代数复算，初始 $3.5$ MeV 处电子减速率约 $2.56\,\mathrm{s^{-1}}$（时间 $0.39$ s），氘约 $0.066\,\mathrm{s^{-1}}$（$15.2$ s），氚约 $0.051\,\mathrm{s^{-1}}$（$19.5$ s）；电子显然首先主导。由 $\nu_{s,i}\propto E^{-3/2}$，到数十至百 keV 后离子份额上升；官方扫描在更低速排序下给氘、氚的局部减速时间约 $0.016$、$0.014$ s。全过程时间须积分 $\int\dd E/[2E\sum_b\nu_{s,b}(E)]$，不能把三个起始局部时间相加或当作常数。上述数字是按官方公式条件重算的对照，具体热化时间仍依末能量和碰撞对数。
\end{solution}

\begin{exercise}\label{ex:hist02-3-3}
PS3 原题3（有官方答案）。均匀环形氢等离子体 $T_e=1\,$keV、$a=0.30\,$m、$R=1.20\,$m、$I=4\times10^5\,$A；取 $\ln\Lambda=16$，在最低阶 $a/R$ 求环向驱动场及一匝 loop voltage。若 $n=10^{13}\,\mathrm{cm^{-3}}$，用电子—离子碰撞估计忽略相对论的逃逸电子最小平行能，再检查相对论假设。
\end{exercise}
\begin{solution}
编者独立解答。Spitzer 电阻率近似 $\eta\simeq5.2\times10^{-5}\ln\Lambda/T_{e,\rm eV}^{3/2}\,\Omega\mathrm m$，因此 $J=I/(\pi a^2)=1.415\times10^6\,$A/m$^2$，$\eta\simeq2.63\times10^{-8}\,\Omega\mathrm m$，$E_\phi=\eta J\simeq3.72\times10^{-2}\,$V/m，$V_{\rm loop}=2\pi RE_\phi\simeq0.281\,$V。高速电子离子拖曳采用 $F_{ei}\simeq ne^4\ln\Lambda/(4\pi\epsilon_0^2m_ev^2)$；令 $eE_\phi=F_{ei}$ 并用 $W=m_ev^2/2$ 得
\begin{equation}W_{\rm run}\simeq\frac{ne^3\ln\Lambda}{8\pi\epsilon_0^2E_\phi}\simeq5.6\times10^4\,\mathrm{eV}.
\end{equation}
若连电子碰撞拖曳也计入，阈值上升一个量级一倍数。此模型的 $W/(m_ec^2)\sim0.11$，非相对论可作初估，但 $v/c\sim0.47$，精细阈值应改用相对论碰撞算子。不能把这里“只用电子—离子”的阈值同 2003 年不同温度且同时纳入电子阻力的题直接去重。
\end{solution}

\originalsection{2002 PS4：绝热闭合与箍缩}
\begin{exercise}\label{ex:hist02-4-1}
PS4 原题1（有官方答案）。从无碰撞 kinetic equation 的总能量矩，在各向同性标量压强、无热流的闭合下得到 $\partial_t(\rho V^2/2+3p/2)+\nabla\cdot[(\rho V^2/2+5p/2)\mathbf V]-qn\mathbf E\cdot\mathbf V=0$；再减去动量式与 $\mathbf V$ 的点积，推出 $D(p/\rho^{5/3})/Dt=0$。
\end{exercise}
\begin{solution}
编者独立推导。Boltzmann 式乘 $mv^2/2$ 并在速度空间积分，使用 $\mathbf v=\mathbf V+\mathbf c$、$\int\mathbf c f=0$、$\int m c_ic_jf=p\delta_{ij}$ 与热流 $\int mc^2\mathbf c f/2=0$，得到题给总能式。将 $\rho D\mathbf V/Dt=qn(\mathbf E+\mathbf V\times\mathbf B)-\nabla p$ 点乘 $\mathbf V$，用连续式 $D\rho/Dt=-\rho\nabla\cdot\mathbf V$，得动能式 $\partial_t(\rho V^2/2)+\nabla\cdot[(\rho V^2/2)\mathbf V]=qn\mathbf E\cdot\mathbf V-\mathbf V\cdot\nabla p$。相减并整理即 $\frac32Dp/Dt+\frac52p\nabla\cdot\mathbf V=0$，故 $D\ln p/Dt=(5/3)D\ln\rho/Dt$。有热流、各向异性或碰撞源时不再是此简单闭合。
\end{solution}

\begin{exercise}\label{ex:hist02-4-2}
PS4 原题2（有官方答案）。给单流体动量式 $\rho D\mathbf v/Dt=\mathbf J\times\mathbf B-\nabla P$，以及含 $\eta\mathbf J$、Hall／电子压力项和电流惯性项的广义 Ohm 式。说明理想 MHD 忽略哪些项；对典型 tokamak 用无量纲排序判断这些忽略是否合理。原题没有固定装置数值，须申明估计值。
\end{exercise}
\begin{solution}
编者独立解答。理想 Ohm 式是 $\mathbf E+\mathbf v\times\mathbf B=0$，所以丢去电阻 $\eta\mathbf J$、Hall $\mathbf J\times\mathbf B/(ne)$、电子压强 $-\nabla p_e/(ne)$ 与电子惯性；在动量方程还已略黏性等闭合项。以尺度 $L$、速度 $V$，相对 $VB$ 的尺度分别约为 $R_m^{-1}=\eta/(\mu_0LV)$，$d_i/L$ 乘适当 $v_A/V$，$\beta_e d_i/(2L)$ 乘同类因子，以及 $(d_e/L)^2$ 乘动态因子。举例取 $n=10^{20}\,\mathrm{m^{-3}}$、$B=5\,$T、$L=0.5\,$m、$T_e=2\,$keV、$V\sim v_A$、氢离子，则 $d_i=c/\omega_{pi}\simeq2.3\,$cm，$d_e\simeq0.53\,$mm，Hall 在全局尺度约 $0.046$，电子惯性几百万分量级；$\beta_e\simeq3.2\times10^{-3}$，压强项在全局尺度也小。热电阻率约 $10^{-8}\,\Omega\mathrm m$，$R_m\gg1$。在重联电流片 $L\sim d_i$、高频或边界层时这些小量不再小，理想 MHD 不适用。该数字是明确指定的示例，不伪称原题给定装置参数。
\end{solution}

\begin{exercise}\label{ex:hist02-4-3}
PS4 原题3（有官方答案）。静态圆柱 Z 箍缩电流 $J_z(r)=c_1(r^2/a^2)/(1+r^2/a^2)^3$，求 $B_\theta(r)$、$p(r)$ 并以总电流 $I$ 表示，描绘剖面。既然电流和压力在远处衰减，这是否违反不能仅靠自身电流有限域约束的 virial 定理？公开作业给出完整题面并注明改编自 J. P. Freidberg 的教材 Problem 5.2；本卷只按课程公开题面求解，未复制教材。
\end{exercise}
\begin{solution}
编者独立解答。设 $x=r^2/a^2$，$I=2\pi\int_0^\infty J_zr\dd r=\pi c_1a^2/2$，即 $c_1=2I/(\pi a^2)$。Amp\`ere 定律和径向平衡 $p'=-J_zB_\theta$ 给
\begin{equation}B_\theta(r)=\frac{\mu_0I}{2\pi r}\frac{x^2}{(1+x)^2},\qquad
p(r)=\frac{\mu_0c_1^2a^2}{8}\left[\frac{u^2}{2}-\frac{2u^3}{3}+\frac{u^4}{4}\right],\quad u=\frac1{1+x},\quad p(\infty)=0.
\end{equation}
近轴 $J_z\propto r^2,B_\theta\propto r^3$，远处 $J_z\propto r^{-4},B_\theta\propto r^{-1},p\propto r^{-4}$。磁场虽然随半径趋零，但总电流非零使磁能每单位轴长 $\int B_\theta^2r\dd r\sim\int\dd r/r$ 对无穷外域对数发散；题设不满足有限能量、无外部回流／边界的 virial 假设。真实回流或外边界承担约束，不能以局部 $p(r)>0$ 反证全局定理。
\end{solution}
